\begin{alist}
\item Καθώς υπάρχουν λογάριθμοι μόνο θετικών αριθμών, απαιτούμε $10^x - 1 > 0$, άρα
\[10^x > 1 \Leftrightarrow 10^x > 10^0 \Leftrightarrow x > 0\]
αφού η συνάρτηση $g(x) = 10^x$ είναι γνησίως αύξουσα στο $\mathbb{R}$.

\item Πρέπει $f(x) > 0 \Leftrightarrow \log(10^x - 1) > 0 \Leftrightarrow \log(10^x - 1) > \log 1$ και καθώς η συνάρτηση $h(x) = \log x$ είναι γνησίως αύξουσα στο $(0, +\infty)$ παίρνουμε $10^x - 1 > 1 \Leftrightarrow 10^x > 2$ σχέση που γράφεται $10^x > 10^{\log 2}$. Ώστε $x > \log 2$, δηλαδή $x \in (\log 2, +\infty)$.

\item Έχουμε $f(x) + x = \log(10^x - 1) + \log 10^x = \log[10^x(10^x - 1)] = \log(10^x \cdot 10^x - 10^x) = \log(10^{2x} - 10^x)$.

\item Πρέπει να λύσουμε την εξίσωση $f(x) = -x \Leftrightarrow f(x) + x = 0 \Leftrightarrow \log(10^{2x} - 10^x) = 0$ άρα $10^{2x} - 10^x = 1 \Leftrightarrow (10^x)^2 - 10^x - 1 = 0$. Θέτοντας $10^x = y > 0$, παίρνουμε την δευτεροβάθμια εξίσωση $y^2 - y - 1 = 0$ με διακρίνουσα $\Delta = (-1)^2 - 4 \cdot 1 \cdot (-1) = 5$, οπότε $y = \dfrac{-(-1) \pm \sqrt{5}}{2} = \dfrac{1 \pm \sqrt{5}}{2}$. Έτσι $y = \dfrac{\sqrt{5} + 1}{2}$, αφού $y > 0$.
Τελικά $10^x = \dfrac{\sqrt{5} + 1}{2} \Leftrightarrow x = \log\left(\dfrac{\sqrt{5} + 1}{2}\right)$.

Άρα το ζητούμενο σημείο είναι το $\left(\log\left(\dfrac{\sqrt{5} + 1}{2}\right), -\log\left(\dfrac{\sqrt{5} + 1}{2}\right)\right)$.
\end{alist}
